\documentclass[12pt]{article}
\usepackage[utf8]{inputenc}
\usepackage[T1]{fontenc}
\usepackage[margin=1in]{geometry}
\usepackage{setspace}
\usepackage{xcolor}
\usepackage[colorlinks=true, urlcolor=blue!60!black, linkcolor=blue!60!black]{hyperref}
\setlength{\parindent}{2em}
\setlength{\parskip}{1.25em}
\renewcommand{\baselinestretch}{1.0}
\title{\textit{United States v. Brown}, 381 U.S. 437 (1965)}
\author{Case Brief}
\date{}
\begin{document}
\maketitle

\section{Case Metadata}

\begin{itemize}
  \item \textbf{Case Name:} \textit{United States v. Brown}, 381 U.S. 437 (1965) (85 S. Ct. 1707; 14 L. Ed. 2d 484).
  \item \textbf{Briefing Book Placement:} Chapter 3: Federal Legislative Power (Article I, \S~9, cl. 3 --- Bill of Attainder Clause).
  \item \textbf{Dates:} Argued March 29, 1965; Decided June 7, 1965.
  \item \textbf{Lower Court:} Certiorari to the U.S. Court of Appeals for the Ninth Circuit (\textit{Brown v. United States}, 334 F.2d 488 (9th Cir. 1964) (en banc), reversing conviction in the U.S. District Court for the Northern District of California).
  \item \textbf{Vote:} 5--4, affirming the Ninth Circuit's judgment dismissing the indictment.
  \item \textbf{Majority Opinion:} Chief Justice Earl Warren, joined by Justices Hugo L. Black, William O. Douglas, William J. Brennan, Jr., and Arthur J. Goldberg.
  \item \textbf{Concurring Opinions:} None.
  \item \textbf{Dissenting Opinion:} Justice Byron R. White, joined by Justices Tom C. Clark, John M. Harlan II, and Potter Stewart.
  \item \textbf{Primary Opinion Sources:} \href{https://supreme.justia.com/cases/federal/us/381/437/}{Justia U.S. Supreme Court Center} and \href{https://www.law.cornell.edu/supremecourt/text/381/437}{Cornell Legal Information Institute (LII)}.
\end{itemize}

\section{Facts and Procedural History}

In 1959, Congress enacted the Labor-Management Reporting and Disclosure Act (LMRDA) to curb union corruption and protect interstate commerce from ``political strikes.'' Section 504 of the LMRDA (29 U.S.C. \S~504) replaced \S~9(h) of the Taft-Hartley Act, which had previously denied National Labor Relations Board (NLRB) access to unions whose officers refused to file non-Communist affidavits. Rather than using an affidavit condition, \S~504 made it a federal crime---punishable by a fine of up to \$10,000, imprisonment for up to one year, or both---for any person who was a member of the Communist Party, or had been a member during the preceding five years, to serve willfully as an officer, executive board member, or non-clerical/non-custodial employee of a labor organization (\href{https://supreme.justia.com/cases/federal/us/381/437/}{381 U.S. at 438--40 \& n.1}).

Respondent Archie Brown had worked as a San Francisco longshoreman and had been an open, avowed member of the Communist Party for more than twenty-five years. Members of Local 10 of the International Longshoremen's and Warehousemen's Union elected Brown to the union's Executive Board for consecutive one-year terms in 1959, 1960, and 1961. On May 24, 1961, a federal grand jury in the Northern District of California indicted Brown on one count of willfully violating \S~504. The government neither alleged nor proved that Brown had ever advocated illegal union conduct or instigated a political strike; indeed, the trial court struck defense evidence that the executive board had never called a strike and rejected an offer to prove that the union had not struck since 1948 (\href{https://supreme.justia.com/cases/federal/us/381/437/}{381 U.S. at 440 \& n.4}).

A jury convicted Brown, and the district court sentenced him to six months' imprisonment. Sitting en banc, the Ninth Circuit reversed and remanded with instructions to dismiss the indictment, holding that \S~504 violated the First and Fifth Amendments (334 F.2d 488). The Supreme Court granted certiorari (379 U.S. 899) and affirmed on an alternative ground raised in Brown's motion to dismiss: that \S~504 was an unconstitutional bill of attainder under Article I, \S~9, cl. 3. Having invalidated the statute on attainder grounds, the Court declined to reach the First and Fifth Amendment issues (\href{https://www.law.cornell.edu/supremecourt/text/381/437}{381 U.S. at 440}).

\section{Issue(s) Presented}

Does \S~504 of the Labor-Management Reporting and Disclosure Act of 1959, which criminalizes service as a labor union officer or non-clerical employee by any current or recent (within five years) member of the Communist Party, violate the federal Bill of Attainder Clause of Article I, \S~9, cl. 3 of the United States Constitution?

\section{Holding}

Yes. In a 5--4 decision written by Chief Justice Warren, the Supreme Court held that \S~504 constitutes an unconstitutional bill of attainder. The Bill of Attainder Clause forbids Congress from legislatively singling out an easily ascertainable group---here, members of the Communist Party---and inflicting punishment upon them without a judicial trial, even when Congress's objective is preventive (forestalling future political strikes) rather than retributive (\href{https://supreme.justia.com/cases/federal/us/381/437/}{381 U.S. at 440, 458--62}).

\section{Reasoning}

Chief Justice Warren structured the Court's opinion around four pillars: constitutional history and separation of powers, the distinction between general rulemaking and legislative specification, the rejection of ``shorthand'' group proxies, and the inclusion of preventive sanctions within ``punishment.''

\begin{enumerate}
  \item \textbf{Separation of Powers and ``Trial by Legislature'':} Adopted unanimously without debate at the Constitutional Convention, the Bill of Attainder Clauses (Art. I, \S~9, cl. 3 and Art. I, \S~10) were not narrow technical bans limited to English parliamentary death sentences with corruption of blood. Drawing on Madison (\textit{Federalist} Nos. 44, 47, 48), Hamilton (\textit{Federalist} No. 78), and Chief Justice Marshall in \textit{Fletcher v. Peck} (1810), Warren framed the Clause as a vital separation-of-powers bulwark confining Congress to general rulemaking while reserving individualized adjudication of blameworthiness and punishment to politically independent courts and juries (\href{https://supreme.justia.com/cases/federal/us/381/437/}{381 U.S. at 441--46}). Precedents such as \textit{Cummings v. Missouri} (1867), \textit{Ex parte Garland} (1867), and \href{https://supreme.justia.com/cases/federal/us/328/303/}{\textit{United States v. Lovett}}, 328 U.S. 303 (1946), confirmed that legislative disqualification from a vocation or office counts as attainder whether targets are named individually or identified as a group.

  \item \textbf{Specification Versus General Rulemaking:} Congress undoubtedly holds Commerce Clause power to bar individuals likely to incite political strikes from union leadership. Yet \S~504 failed to articulate general behavioral criteria for courts or juries to apply. Instead, Congress itself designated Communist Party members as possessing the feared traits. Warren contrasted \S~504 with the Subversive Activities Control Act upheld in \textit{Communist Party v. Subversive Activities Control Board}, 367 U.S. 1 (1961), which regulated defined activities rather than an enumerated organization and required an administrative hearing with judicial review (\href{https://www.law.cornell.edu/supremecourt/text/381/437}{381 U.S. at 449--52}).

  \item \textbf{Distinguishing Conflict-of-Interest Laws and Rejecting ``Shorthand'':} The Solicitor General analogized \S~504 to \S~32 of the Banking Act of 1933 (upheld in \textit{Board of Governors v. Agnew}, 329 U.S. 441 (1947)), which barred securities underwriters from serving as bank directors. Warren rejected the analogy on three grounds: (a) \S~504 targeted a political group, the classic historical target of attainders; (b) \S~32 reflected general human psychology applicable to anyone in a dual financial role, rather than censuring a specific faction; and (c) \S~32 established an objective conduct standard with an administrative escape clause. Party membership could not serve as a neutral ``shorthand'' for strike-incitement propensities because political beliefs are personal and not all Party members share illegal objectives (\textit{Schneiderman v. United States}; \textit{Aptheker v. Secretary of State}; \textit{Schware v. Board of Bar Examiners}) (\href{https://supreme.justia.com/cases/federal/us/381/437/}{381 U.S. at 453--56}).

  \item \textbf{Preventive Sanctions Are Still ``Punishment'':} The government relied on \href{https://supreme.justia.com/cases/federal/us/339/382/}{\textit{American Communications Ass'n v. Douds}}, 339 U.S. 382 (1950), which had upheld Taft-Hartley \S~9(h) as preventive rather than retributive because officers could resign from the Party to escape disqualification. Warren noted that \S~504's five-year lookback eliminated immediate escapability, but went further to repudiate \textit{Douds}'s premise: ``punishment'' encompasses retributive, rehabilitative, deterrent, and preventive aims alike. Historical English and Revolutionary-era attainders against Tories were frequently preventive, and \textit{Douds} misread \href{https://supreme.justia.com/cases/federal/us/328/303/}{\textit{Lovett}}, which had also sought to purge ``subversives'' prospectively (\href{https://supreme.justia.com/cases/federal/us/381/437/}{381 U.S. at 456--61}).
\end{enumerate}

\section{Concurring and Dissenting Opinions}

No separate concurring opinions were filed.

Justice White dissented, joined by Justices Clark, Harlan, and Stewart (\href{https://supreme.justia.com/cases/federal/us/381/437/}{381 U.S. at 462--78}). White argued that the majority abandoned the Court's established multi-factor framework for distinguishing punishment from regulation (\textit{Flemming v. Nestor}; \textit{Kennedy v. Mendoza-Martinez}) in favor of a rigid per se rule equating legislative specificity with judicial usurpation. In White's view, the majority's twin objections---that \S~504 was simultaneously underinclusive (omitting non-Communists who might call political strikes) and overinclusive (reaching benign Party members)---threatened well-settled regulatory statutes, including \href{https://supreme.justia.com/cases/federal/us/339/382/}{\textit{Douds}} (which White observed was effectively overruled), conflict-of-interest laws (\textit{Agnew}), felon-disqualification rules (\textit{Hawker v. New York}; \textit{De Veau v. Braisted}), and national-security employment bars.

White further questioned the majority's premise that the Bill of Attainder Clause enshrines a strict federal or state separation of powers, pointing to private bills, legislative contempt, and the Constitution's refusal to mandate rigid separation of powers for the states. Legislatures routinely investigate concrete abuses and regulate specific groups piecemeal (\textit{Patsone v. Pennsylvania}; \textit{Railway Express Agency v. New York}). Applying the \textit{Mendoza-Martinez} factors, White concluded that \S~504 was rationally tied to the legitimate, nonpunitive goal of preventing political strikes and rested on prior congressional and Subversive Activities Control Board findings rather than punitive animus (\href{https://www.law.cornell.edu/supremecourt/text/381/437}{381 U.S. at 472--78}).

\section{Key Quotations}

\begin{quote}
``The best available evidence, the writings of the architects of our constitutional system, indicates that the Bill of Attainder Clause was intended not as a narrow, technical (and therefore soon to be outmoded) prohibition, but rather as an implementation of the separation of powers, a general safeguard against legislative exercise of the judicial function, or, more simply, trial by legislature.'' --- Chief Justice Warren, \href{https://supreme.justia.com/cases/federal/us/381/437/}{381 U.S. at 442}.
\end{quote}

\begin{quote}
``It would be archaic to limit the definition of `punishment' to `retribution.' Punishment serves several purposes; retributive, rehabilitative, deterrent---and preventive. One of the reasons society imprisons those convicted of crimes is to keep them from inflicting future harm, but that does not make imprisonment any the less punishment.'' --- Chief Justice Warren, \href{https://supreme.justia.com/cases/federal/us/381/437/}{381 U.S. at 458}.
\end{quote}

\begin{quote}
``We do not hold today that Congress cannot weed dangerous persons out of the labor movement, any more than the Court held in \textit{Lovett} that subversives must be permitted to hold sensitive government positions. Rather, we make again the point made in \textit{Lovett}: that Congress must accomplish such results by rules of general applicability. It cannot specify the people upon whom the sanction it prescribes is to be levied. Under our Constitution, Congress possesses full legislative authority, but the task of adjudication must be left to other tribunals.'' --- Chief Justice Warren, \href{https://supreme.justia.com/cases/federal/us/381/437/}{381 U.S. at 461}.
\end{quote}

\begin{quote}
``In the Court's view, therefore, \S~504 is too narrow in specifying the particular class; but it is also too broad in treating all members of the class alike. On both counts---underinclusiveness and overinclusiveness---\S~504 is invalid as a bill of attainder because Congress has engaged in forbidden factfinding about individuals and groups, and has thus strayed into the area reserved to the judiciary by the Constitution.'' --- Justice White, dissenting, \href{https://supreme.justia.com/cases/federal/us/381/437/}{381 U.S. at 464}.
\end{quote}

\section{Doctrinal Significance and Modern Legacy (Through 2026)}

\textit{United States v. Brown} marks the high-water mark of the Supreme Court's expansive, separation-of-powers interpretation of the Bill of Attainder Clause---and remains, as of 2026, the last occasion on which the Supreme Court invalidated a federal statute as an unconstitutional bill of attainder. Within a dozen years, however, the Court pivoted substantially toward Justice White's functionalist dissent.

In \href{https://supreme.justia.com/cases/federal/us/433/425/}{\textit{Nixon v. Administrator of General Services}}, 433 U.S. 425 (1977), the Court upheld the Presidential Recordings and Materials Preservation Act, which named former President Richard Nixon specifically, holding that Nixon constituted a ``legitimate class of one'' and that specificity alone does not create an attainder. \textit{Nixon} established the modern three-part inquiry for ``punishment'': (1) the historical test (whether the burden falls within traditional attainder penalties), (2) the functional test (whether the law reasonably furthers nonpunitive legislative purposes), and (3) the motivational test (whether the legislative record evinces punitive intent). Seven years later, \href{https://supreme.justia.com/cases/federal/us/468/841/}{\textit{Selective Service System v. Minnesota Public Interest Research Group}}, 468 U.S. 841 (1984), reinforced \textit{Nixon} by upholding a law denying federal student aid to non-registrants for the draft, emphasizing both its nonpunitive compliance goal and students' ability to qualify by registering.

Through 2026, \textit{Nixon}'s narrowing of \textit{Brown} has proven decisive in major national-security and technology-sanctions challenges where Congress targeted named foreign-linked corporations. Courts rejected attainder challenges to National Defense Authorization Act procurement bans naming Russian and Chinese firms in \textit{Kaspersky Lab, Inc. v. DHS}, 909 F.3d 446 (D.C. Cir. 2018) (summarized in \href{https://www.everycrsreport.com/files/2026-09-02_IF11293_c333cd4fdeb9bdf460eadc723c3b9ea6c1a30e27.html}{CRS Report IF11293 (2026)}), and \href{https://law.justia.com/cases/federal/district-courts/texas/txedce/4:2019cv00159/188186/51/}{\textit{Huawei Technologies USA, Inc. v. United States}}, 440 F. Supp. 3d 607 (E.D. Tex. 2020), reasoning that prospective cybersecurity restrictions are regulatory rather than punitive. Most prominently, in \href{https://law.justia.com/cases/federal/appellate-courts/cadc/24-1113/24-1113-2024-12-06.html}{\textit{TikTok Inc. v. Garland}}, 122 F.4th 930 (D.C. Cir. 2024), the D.C. Circuit unanimously rejected ByteDance and TikTok's Bill of Attainder challenge to the Protecting Americans from Foreign Adversary Controlled Applications Act (PAFACAA), holding that even though the statute named TikTok specifically, compelled divestiture was a preventive national-security measure rather than historical confiscation or punishment. When the Supreme Court affirmed in \href{https://supreme.justia.com/cases/federal/us/604/24-656/}{\textit{TikTok Inc. v. Garland}}, 604 U.S. \_\_\_ (2025) (per curiam), it granted review solely on the First Amendment question, leaving the D.C. Circuit's narrow reading of \textit{Brown} intact. State appellate decisions through 2025--2026 likewise confirm the steep post-\textit{Nixon} bar for attainder plaintiffs (e.g., \href{https://law.justia.com/cases/california/court-of-appeal/2025/a170255.html}{\textit{Hastings College Conservation Committee v. State}} (Cal. Ct. App. 2025)), leaving \textit{Brown} most potent when Congress imposes criminal or vocational disqualifications directly on natural persons defined by political association.

\section{Pedagogical Critiques and Exam Hypotheticals}

\subsection*{Key Critiques}

\begin{itemize}
  \item \textbf{Doctrinal Instability After \textit{Nixon}:} \textit{Brown} insists that preventive disqualifications targeting an ascertainable group constitute punishment, yet \href{https://supreme.justia.com/cases/federal/us/433/425/}{\textit{Nixon}} and modern circuit decisions treat preventive national-security targeting as nonpunitive regulation. Students should ask whether \textit{Brown} survives mainly as a rule against criminalizing political-party status rather than a general ban on targeted prophylaxis.
  \item \textbf{First Amendment Shadow:} The Ninth Circuit struck down \S~504 under the First and Fifth Amendments, and Warren's ``shorthand'' analysis leans heavily on First Amendment overbreadth cases (\textit{Aptheker}, \textit{Schware}). Critics note that \textit{Brown} may use Article I's structural Attainder Clause to do First Amendment associational-freedom work while sidestepping messy scrutiny tiers.
\end{itemize}

\subsection*{Exam Hypothetical and ``Opposite Side'' Contrast}

\textbf{Hypothetical (Attainder Likely Under \textit{Brown}):} After congressional hearings link an activist coalition, ``GridLock Now,'' to disruptive energy protests, Congress enacts a federal statute making it a felony for any current member---or anyone who belonged to GridLock Now during the past three years---to serve as an officer or board member of any interstate electric utility. A utility director who joined GridLock Now two years ago to support solar incentives, and never advocated disruption, is indicted. \textit{Analysis:} Under \textit{Brown}, the statute fails as a bill of attainder: it singles out an ascertainable political group by name, imposes a historical vocational bar backed by criminal punishment and a multi-year lookback, and substitutes legislative factfinding about group propensities for individualized judicial adjudication, regardless of Congress's preventive grid-security rationale.

\textbf{``Opposite Side'' Contrast (Permissible General Rulemaking / Regulation):} Contrast two statutes that would survive attainder scrutiny. First, a statute barring any person from serving as an electric-utility officer if a federal agency or court finds, after an individualized hearing with judicial review, that the person conspired to sabotage critical infrastructure mirrors the general behavioral rulemaking Warren endorsed in \textit{Brown} (\href{https://supreme.justia.com/cases/federal/us/381/437/}{381 U.S. at 461}) and \textit{Communist Party v. SACB}. Second, a statute prohibiting federal procurement of grid-control software from a named foreign-adversary-controlled vendor (or requiring foreign divestiture) would likely be upheld under \href{https://supreme.justia.com/cases/federal/us/433/425/}{\textit{Nixon}}, \href{https://law.justia.com/cases/federal/district-courts/texas/txedce/4:2019cv00159/188186/51/}{\textit{Huawei}}, and \href{https://law.justia.com/cases/federal/appellate-courts/cadc/24-1113/24-1113-2024-12-06.html}{\textit{TikTok}} as nonpunitive, prospective national-security regulation rather than punitive disqualification of natural persons for political association.

\end{document}
